\section{Design Philosophy}

VOO's fundamental design decision is that \textbf{objects are values, not references}.
A VOO object is a plain Tcl list---creating a copy with \texttt{set p2 \$p1} shares the
reference via copy-on-write, and modifying \texttt{p2} leaves \texttt{p1} unchanged.
This eliminates destructor management entirely: objects are garbage-collected when their
reference count reaches zero, preventing the memory leak vulnerabilities inherent in
manual destruction.

Classes are implemented as standard Tcl namespaces. Field indices are stored as
namespace variables, and accessors are ordinary procedures. No specialized infrastructure
is required. VOO constructs OO capabilities from established patterns: lists for object
storage, dictionaries for map fields, namespaces for encapsulation, procedures for
methods, \texttt{upvar} for setter-by-reference, and Tcl's COW for automatic memory
management. This yields seamless integration with existing codebases, predictable
performance from well-understood primitives, and simplified debugging through direct
string representation.

Performance advantages follow directly from this simplicity. Getters are a single
\texttt{lindex} with a precomputed index; setters use \texttt{upvar} and \texttt{lset}
for in-place modification; objects carry no wrapper overhead. Copy-on-write is delegated
entirely to Tcl's existing mechanisms.

\begin{lstlisting}[language=Tcl, caption={VOO getter and setter --- the entire runtime implementation}]
# VOO getter and setter --- the entire runtime implementation
proc Point::get.x {this} { return [lindex $this 0] }
proc Point::set.x {thisVar value} { upvar $thisVar this; lset this 0 $value }
\end{lstlisting}
